% GENERATED FILE. Edit canonical lessons, then run npm run generate:books.
% canonical-source-hash: fnv1a64:263e0a1f785389d9
% canonical-lessons: GU-C455-gamgin, GU-C455-vyatha, GU-C455-nisaso, GU-C455-radavvu, GU-C455-rudan

\chapter{A Heavy Heart}
\label{ch:gu-a-heavy-heart}
\hlchaptermodality{455}

\begin{chapteropening}
\textbf{By the end of this chapter:} \emph{I can say how grief shows on a person.}

Complete the last lesson of chapter 455: I can say how grief shows on a person.
\end{chapteropening}

\section[\texorpdfstring{\gu{ગમગીન} (gamgīn)}{ગમગીન (gamgīn)}]{\gu{ગમગીન} (gamgīn) — gloomy; sorrowful}

\label{lesson:GU-C455-gamgin}



\begin{quote}
Before the new one: say the Gujarati for distraught, then the Gujarati for to fidget.
\end{quote}

\subsection*{You'll want to know: \gu{ગમગીન}}
\textbf{\gu{ગમગીન}} — \emph{gamgīn} — \textquotedblleft{}gloomy; sorrowful\textquotedblright{}.

\begin{quote}
\gu{મિત્ર} \gu{ગયા} \gu{પછી} \gu{અરુણ} \gu{આખો} \gu{દિવસ} \gu{ગમગીન} \gu{રહ્યો}.
\end{quote}

\begin{quote}
\emph{mitra gayā pachī Aruṇ ākho divas gamgīn rahyo.}
\end{quote}

\begin{quote}
After his friend left, Arun stayed gloomy all day.
\end{quote}

\begin{cousinweb}
\textbf{\gu{ગમગીન}} (\emph{gamgīn}) is usually traced to Persian \emph{ġamgīn}, \textquotedblleft{}sorrowful\textquotedblright{}, built on \emph{ġam}, \textquotedblleft{}grief\textquotedblright{}.
\end{cousinweb}

\begin{grammarlens}[title={What you've built}]
One more word for describing what you lived and felt.
\end{grammarlens}

\subsection*{Your turn}
\begin{itemize}
\raggedright
  \item \emph{Say it:} \emph{gamgīn}
  \item \emph{Say it:} \emph{gamgīn}, once more
  \item \emph{Say it:} \emph{mitra gayā pachī Aruṇ ākho divas gamgīn rahyo.} — After his friend left, Arun stayed gloomy all day.
  \item \emph{Recall:} say the Gujarati for distraught, then the Gujarati for to fidget, then say \emph{gamgīn} again
\end{itemize}

\begin{tcolorbox}[breakable,colback=teal!4,colframe=teal!35!black,title={Before you move on}]
What does \gu{ગમગીન} mean? (\textquotedblleft{}Gloomy; sorrowful\textquotedblright{}.) Say it once more.
\end{tcolorbox}
\section[\texorpdfstring{\gu{વ્યથા} (vyathā)}{વ્યથા (vyathā)}]{\gu{વ્યથા} (vyathā) — anguish}

\label{lesson:GU-C455-vyatha}



\begin{quote}
Before the new one: say the Gujarati for to fidget, then the Gujarati for gloomy.
\end{quote}

\subsection*{You'll want to know: \gu{વ્યથા}}
\textbf{\gu{વ્યથા}} — \emph{vyathā} — \textquotedblleft{}anguish\textquotedblright{}.

\begin{quote}
\gu{તેણે} \gu{પોતાની} \gu{વ્યથા} \gu{કોઈને} \gu{કહી} \gu{નહીં}.
\end{quote}

\begin{quote}
\emph{teṇe potānī vyathā koīne kahī nahī̃.}
\end{quote}

\begin{quote}
She told no one about her anguish.
\end{quote}

A feminine noun.

\begin{cousinweb}
\textbf{\gu{વ્યથા}} (\emph{vyathā}) is a tatsama, Sanskrit \emph{vyathā}, \textquotedblleft{}agitation, uneasiness, pain, anguish\textquotedblright{}, as Monier-Williams records it.
\end{cousinweb}

\begin{grammarlens}[title={What you've built}]
One more word for describing what you lived and felt.
\end{grammarlens}

\subsection*{Your turn}
\begin{itemize}
\raggedright
  \item \emph{Say it:} \emph{vyathā}
  \item \emph{Say it:} \emph{vyathā}, once more
  \item \emph{Say it:} \emph{teṇe potānī vyathā koīne kahī nahī̃.} — She told no one about her anguish.
  \item \emph{Recall:} say the Gujarati for to fidget, then the Gujarati for gloomy, then say \emph{vyathā} again
\end{itemize}

\begin{tcolorbox}[breakable,colback=teal!4,colframe=teal!35!black,title={Before you move on}]
What does \gu{વ્યથા} mean? (\textquotedblleft{}Anguish\textquotedblright{}.) Say it once more.
\end{tcolorbox}
\section[\texorpdfstring{\gu{નિસાસો} (nisāso)}{નિસાસો (nisāso)}]{\gu{નિસાસો} (nisāso) — a sigh}

\label{lesson:GU-C455-nisaso}



\begin{quote}
Before the new one: say the Gujarati for gloomy, then the Gujarati for anguish.
\end{quote}

\subsection*{You'll want to know: \gu{નિસાસો}}
\textbf{\gu{નિસાસો}} — \emph{nisāso} — \textquotedblleft{}a sigh\textquotedblright{}.

\begin{quote}
\gu{દાદીએ} \gu{જૂના} \gu{ફોટા} \gu{જોઈને} \gu{નિસાસો} \gu{નાખ્યો}.
\end{quote}

\begin{quote}
\emph{dādīe jūnā phoṭā joīne nisāso nākhyo.}
\end{quote}

\begin{quote}
Looking at the old photos, my grandmother sighed.
\end{quote}

A masculine noun.

\begin{cousinweb}
\textbf{\gu{નિસાસો}} (\emph{nisāso}) is inherited from Sanskrit \emph{niḥśvāsa}, \textquotedblleft{}a sigh\textquotedblright{}, a breathing out; CDIAL lists the Gujarati word (CDIAL 7111).
\end{cousinweb}

\begin{grammarlens}[title={What you've built}]
One more word for describing what you lived and felt.
\end{grammarlens}

\subsection*{Your turn}
\begin{itemize}
\raggedright
  \item \emph{Say it:} \emph{nisāso}
  \item \emph{Say it:} \emph{nisāso}, once more
  \item \emph{Say it:} \emph{dādīe jūnā phoṭā joīne nisāso nākhyo.} — Looking at the old photos, my grandmother sighed.
  \item \emph{Recall:} say the Gujarati for gloomy, then the Gujarati for anguish, then say \emph{nisāso} again
\end{itemize}

\begin{tcolorbox}[breakable,colback=teal!4,colframe=teal!35!black,title={Before you move on}]
What does \gu{નિસાસો} mean? (\textquotedblleft{}A sigh\textquotedblright{}.) Say it once more.
\end{tcolorbox}
\section[\texorpdfstring{\gu{રડાવવું} (raḍāvvũ)}{રડાવવું (raḍāvvũ)}]{\gu{રડાવવું} (raḍāvvũ) — to make someone cry}

\label{lesson:GU-C455-radavvu}



\begin{quote}
Before the new one: say the Gujarati for anguish, then the Gujarati for a sigh.
\end{quote}

\subsection*{You'll want to know: \gu{રડાવવું}}
\textbf{\gu{રડાવવું}} — \emph{raḍāvvũ} — \textquotedblleft{}to make someone cry\textquotedblright{}.

\begin{quote}
\gu{આ} \gu{ફિલ્મ} \gu{બધાને} \gu{રડાવે} \gu{છે}.
\end{quote}

\begin{quote}
\emph{ā philm badhāne raḍāve che.}
\end{quote}

\begin{quote}
This film makes everyone cry.
\end{quote}

Past, with \gu{તેણે} (he or she …): \gu{તેણે} \gu{બાળકને} \gu{રડાવ્યું}, \emph{teṇe bāḷakne raḍāvyũ}.

\begin{cousinweb}
\textbf{\gu{રડાવવું}} (\emph{raḍāvvũ}) is the causative of \textbf{\gu{રડવું}} (\emph{raḍvũ}), \textquotedblleft{}to cry\textquotedblright{}, which is inherited from Sanskrit \emph{raṭati}, \textquotedblleft{}cries, yells\textquotedblright{} (CDIAL 10590).
\end{cousinweb}

\begin{grammarlens}[title={What you've built}]
One more word for describing what you lived and felt.
\end{grammarlens}

\subsection*{Your turn}
\begin{itemize}
\raggedright
  \item \emph{Say it:} \emph{raḍāvvũ}
  \item \emph{Say it:} \emph{raḍāvvũ}, once more
  \item \emph{Say it:} \emph{ā philm badhāne raḍāve che.} — This film makes everyone cry.
  \item \emph{Recall:} say the Gujarati for anguish, then the Gujarati for a sigh, then say \emph{raḍāvvũ} again
\end{itemize}

\begin{tcolorbox}[breakable,colback=teal!4,colframe=teal!35!black,title={Before you move on}]
What does \gu{રડાવવું} mean? (\textquotedblleft{}To make someone cry\textquotedblright{}.) Say it once more.
\end{tcolorbox}
\section[\texorpdfstring{\gu{રુદન} (rudan)}{રુદન (rudan)}]{\gu{રુદન} (rudan) — weeping; wailing}

\label{lesson:GU-C455-rudan}



\begin{quote}
Before the new one: say the Gujarati for a sigh, then the Gujarati for to make someone cry.
\end{quote}

\subsection*{You'll want to know: \gu{રુદન}}
\textbf{\gu{રુદન}} — \emph{rudan} — \textquotedblleft{}weeping; wailing\textquotedblright{}.

\begin{quote}
\gu{બાળકનું} \gu{રુદન} \gu{આખી} \gu{રાત} \gu{ચાલ્યું}.
\end{quote}

\begin{quote}
\emph{bāḷaknũ rudan ākhī rāt cālyũ.}
\end{quote}

\begin{quote}
The child's crying went on all night.
\end{quote}

A neuter noun.

\begin{cousinweb}
\textbf{\gu{રુદન}} (\emph{rudan}) is a tatsama, Sanskrit \emph{rudana}, \textquotedblleft{}the act of crying, weeping, lamentation\textquotedblright{}, as Monier-Williams records it.
\end{cousinweb}

\begin{grammarlens}[title={What you've built}]
One more word for describing what you lived and felt.
\end{grammarlens}

\subsection*{Your turn}
\begin{itemize}
\raggedright
  \item \emph{Say it:} \emph{rudan}
  \item \emph{Say it:} \emph{rudan}, once more
  \item \emph{Say it:} \emph{bāḷaknũ rudan ākhī rāt cālyũ.} — The child's crying went on all night.
  \item \emph{Recall:} say the Gujarati for a sigh, then the Gujarati for to make someone cry, then say \emph{rudan} again
\end{itemize}

\begin{tcolorbox}[breakable,colback=teal!4,colframe=teal!35!black,title={Before you move on}]
What does \gu{રુદન} mean? (\textquotedblleft{}Weeping; wailing\textquotedblright{}.) Say it once more.
\end{tcolorbox}
